\documentclass[10pt]{article}

\usepackage[utf8]{inputenc}

\usepackage[T1]{fontenc}

\usepackage{amsmath}

\usepackage{amsfonts}

\usepackage{amssymb}

\usepackage[version=4]{mhchem}

\usepackage{stmaryrd}

\usepackage{bbold}

\usepackage{graphicx}

\usepackage[export]{adjustbox}

\graphicspath{ {./images/} }



\begin{document}

товой теории во втором порядке по заряду $e$ возникает аномалия, и вместо сохранения тока получаем



\[

\frac{d_{j 5}^{\lambda}}{\partial_{x}^{\lambda}} \sim e^{2} e^{\mu \nu \rho \sigma} F_{\mu \nu} F_{\rho \sigma} .

\]



Вторая теорема Нетер. Помимо обсуждавшейся выше Н.т., к-рую принято называть пе р в ой Н.т., существует вторая Н.т., к-рая касается тождеств, вытекающих из инвариантности действия относительно преобразований, зависящих от непрерывного параметра, то есть от произвольной функции. Наибольшее значение она получает в применении к случаю «полей материи», взаимодействующих с калибровочным полем $A(x)$ - полем, физич. содержание к-рого не меняется при определенных зависящих от произвольной функшии $\lambda(x)$ преобразованиях, называемых преобразованиями калибровки. Вычисляя вариацию действия для поля материи во внешнем калибровочном поле, вызванную бесконечно малым калибровочным преобразованием $\delta \lambda(x), \delta \lambda(x)=0$ на границах области интегрирования, следует учитывать только вызываемые изменением калибровки вариации калибровочного поля $\delta_{\lambda} A=\nabla \delta \lambda$ (здесь $\nabla$ - ковариантная производная), поскольку при вариациях полей материи коэффициентами будут левые части уравнений движения. Поэтому



\[

\delta_{\lambda} S=\int d^{4} x \frac{\partial L}{\partial A} \nabla \delta \lambda=-\int d^{4} x \nabla\left(\frac{\partial L}{\partial A}\right) \delta \lambda,

\]



откуда в силу произвольности $\delta \lambda(x)$ вытекает к о в а р и а н тный закон сохранения:



\[

\nabla j=0 ; \quad j=\frac{\partial L}{\partial A} .

\]



При обращении в формуле (8) внешнего калибровочного поля в нуль ковариантный закон сохранения превращается в обычный, получаемый по первой Н.т. Подчеркнем, что вторая Н.т. приводит к ограничениям на поля материи, исходя из особенностей калибровочного поля. Таким образом, она устанавливает соответствие между свойствами материальных систем и полей, с к-рыми они могут взаимодействовать. Поскольку в правых частях уравнений движения самих калибровочных полей стоят как раз токи (8), то вторая Н. т. налагает тождественные соотношения на левые части этих уравнений. В современной квантовой теории поля вторая Н.т. используется в электродинамике, теории Янга-Миллса полей, гравитации, супергравитации и т. д.



Jum.: [1] Hilbert D., Die Grundlagen der Physik, «Gött. Nachr., Math. Phys. Klasse», 1915, S. 395; [2] K le in F., Die Differentialgesetze fur die Erhaltung von Impuls und Energie in der Einsteinsch. Gravitationstheorie, там же, 1918, S. 171; [3] N other E., Invariante Variationsprobleme, там же, 1918, S. 235; [4] Ландау Л. Д., Л и фш и ц Е. М., Теория поля, 7 изд., М., 1988; [5] П а ул и В., Релятивистская теория элементарных частиц, пер. с англ., М., 1947; [6] Б о го л юбов Н. Н., Ш и рко в Д. В., Введение в теорию квантованных полей, 4 изд., М., 1984; [7] Гельфанд И. М., Фо м ин С. В., Вариационное исчисление, М., 1961; [8] Медведе в Б. В., Начала теоретической физики, М., 1977.



Б. В. Медведев, П. Б. Медведев.



HETEPOB OПEPATOP - см. Интегральное уравнение.



НЕУДЕРХИВАЮЩАЯ СВЯЗЬ - см. Связи механические. НЕУПОРЯДОЧЕННАЯ ФАЗА - см. Дальний $и$ ближний порядок.



НЕУПРУГИЙ ПРОЦЕСС - см. Рассеяния теория.



НЕУСТОЙЧИВОЕ ДВИХЕНИЕ - сМ. Устойчивость движения.



НЕУСТОЙЧИВОЕ МНОГООБРАЗИЕ - см. Центральное многообразие.



НЕУСТОЙЧИВОЕ СЛОЕНИЕ - см. Слоение.



НЕУСТОЙЧИВОСТЬ - см. Устойчивості.



НЕУСТОЙЧИВЫЙ ПОТЕНЦИАЛ ВЗАИМОДЕЙСТВИЯ,



катастрофический потенциал взаимодейст-



в и я,- взаимодействия потенциал, не удовлетворяющий ус- ловию устойчивости (см. Устойчивый потенциал взаимодействия).



Системы частиц с Н. п. в. не имеют термодинамич. поведения (то есть не допускают описания с помощью термодинамич. методов). В системах с Н. П. в. конфигурациям, в K-рых в ограниченных областях пространства $\mathbb{R}^{v}$ собирается много частиц, приписывается большой вес и, следовательно, возможно сосредоточение всех частиц системы в малых областях пространства (коллапс). В силу этого Н. п. в. не используются в статистич. механике.



Пр и м е р 1. Потенциал, принимающий отрицательные значения в нек-рой окрестности начала координат.



П р и м е р 2. Потенциал



\[

\Phi(x)= \begin{cases}x, & 0 \leqslant|x|<R-\delta, \\ -1, & R-\delta \leqslant|x| \leqslant R+\delta, \\ 0, & |x|>R+\delta\end{cases}

\]



является Н. П. в. при $x<12$.



Лит. см. при ст. Устойчивый потенциал взаимодействия. П. В. Малымев.



НЕУСТОЙЧИВЫЙ ПРЕДЕЛЬНЫЙ ЦИКЛ - см. ПредеЛЬный цикл.



НЕХОРОШЕВА TEOPЕМА - теорема, доставляющая эКспоненциальную оценку сверху средней скорости эволюции в возмущенных гамильтоновых системах (см. [2]-[4]). Рacсматриваемая гамильтонова система с аналитич. гамильтонианом



\[

H=H_{0}(I)+\varepsilon H_{1}(I, \varphi, \varepsilon), 0<\varepsilon \ll 1,

\]



где $I=\left(I_{1}, \ldots, I_{n}\right)$ и $\varphi=\left(\varphi_{1}, \ldots, \varphi_{n}\right)-$ сопряженные канонич. переменные, гамильтониан $2 \pi$-периодичен по всем $\varphi_{i}$.



\includegraphics[max width=\textwidth, center]{2023_07_08_2e07cc323f4fca590c1dg-1}

Тогда при достаточно малом $\varepsilon$ для решения $I(t), \varphi(t)$ выполнено



\[

|I(t)-I(0)|<\varepsilon^{b} \text { при } 0 \leqslant t \leqslant \frac{1}{\varepsilon} \exp \left(\frac{1}{\varepsilon^{a}}\right),

\]



если начальное значение $I(0)$ взято внутри $G$. Постоянные $a>0, b>0$ зависят от геометрич. свойств функции $H_{0}$ и оцениваются через значения коэффициентов и показателей крутизны, входящих в определение крутой функции (см. [3]).



Аналогичные оценки справедливы для изменения $I$ в системах с гамильтонианом вида:



\[

H=H_{0}(I)+\varepsilon H_{1}(I, \varphi, p, q),

\]



где $(I, \varphi),(p, q)$ - пары сопряженных канонич. переменных.



Если функция $H_{0}$ - не крутая, то значение $I$ может изменяться на 1 за время $1 / \varepsilon$ (см. [2]). Экспоненциально медленный уход переменной $I$ от начального значения при крутой функции $H_{0}$ называется Арнольда диффузией; в имеющихся примерах скорость этого ухода есть $O(\exp (-\operatorname{const} / \sqrt{\varepsilon}))($ cм. [1]).



Лит.: [1] Арнольд В. И., «Доклады АН СССР», 1964, т. 156, № 1, с. 9-12; [2] Н е хо р о ше в Н. Н., «Функциональный анализ и его прилож.», 1971, т. 5, в. 4, с. 82-83; [3] е го же, «Успехи математич. наук», 1977, т. 32, в. 6, с. 5-66; [4] е го ж е, «Труды семинара им. И. Г. Петровского", 1979, в. 5, с. 5-50. А. И. Нейштадт. НЕЦЕНТРАЛЬНЫЙ УДАР - СМ. УӘар.



НЕЭКВИВАЛЕНТНЫЕ ПРЕДСТАВЛЕНИЯ П е е становочных соотноше ни й - см. Перестановочных соотношений представления неэквивалентные.



НИДИНГ-ИНВАРИАНТ - см. Символическая динамика. НИДИНГ-ПОСЛЕДОВАТЕЛЬНОСТЬ - см. СИмволическая динамика.



НИЗКОЭНЕРГЕТИЧЕСКАЯ TEOРЕМА - См. Токов алгебра.



НИЛЬПОТЕНТНАЯ АЛГЕБРА ЛИ - см. Ли алгебра.



НильПОТЕНтНАЯ ГPУกпА - груnпа $G$, обладающая цепочкой подгрупп



\[

G=G_{0} \supset G_{1} \supset \ldots \supset G_{k}=\{e\}

\]





\end{document}